% PNAS research reports begin directly, without an explicit Introduction
% heading. Keep this file to four principal paragraphs so the first-page column
% transition controlled by \firstpage{4} remains predictable.

Climate-scale deployment of geologic carbon storage requires confidence that
injected \coTwo{} remains contained over operational and post-injection time
scales. Faults are central to that assessment because they can impede flow,
transmit fluids along their planes, or connect otherwise separated units.
Their influence cannot always be reduced to a binary seal-or-conduit label:
the same fault--seal setting can support different migration pathways as its
unresolved internal structure changes. Reliable site assessment therefore
requires knowing not only the expected fault behavior, but also whether
unresolved fault structure can change the containment conclusion.

This requirement creates a multiscale problem. Field characterization
constrains aspects of stratigraphy, fault throw, and burial history, whereas
the local arrangement and continuity of fault-core materials remain
unresolved. That architecture controls anisotropic permeability, porosity,
capillary entry, and phase mobility, while field-scale simulations represent
the fault with much coarser elements. A single representative property field
implicitly assumes that unresolved local variability does not materially
alter the reservoir-scale prediction. That assumption has not been tested
systematically across a broad fault--seal geologic design space.

Stochastic fault-property modeling and multiphase upscaling begin to bridge
this scale gap. PREDICT links local stratigraphy and fault throw to
distributions of three-dimensional anisotropic fault permeability
\cite{SaloSalgado2023}. Pettersson et al.\ upscale stochastic fault properties,
preserve dependence among flow functions, and propagate them to field-scale
leakage distributions, while identifying fully three-dimensional fault
variation as an unresolved limitation \cite{Pettersson2025}. Field-scale
studies further demonstrate that fault representation can influence modeled
\coTwo{} migration and storage-security metrics
\cite{SaloSalgado2025,Lu2025}. What remains unclear is how the importance of
stochastic fault-core structure changes across fault--seal geologies and when
one representative fault gives a misleading reservoir-scale prediction.

Here, we use an end-to-end uncertainty-quantification framework to determine
(i) which characteristics of the fault--seal system control \coTwo{} migration
and containment, (ii) in which geologic settings stochastic fault-core
realizations produce different containment outcomes, and (iii) when a
representative fault is sufficient for storage assessment. We connect
stochastic local three-dimensional fault-core modeling, physics-based
multiphase upscaling, conditionally independent full-fault property sampling,
and field-scale simulation across 162 prescribed geologic scenarios. A
balanced initial ensemble supports comparisons across all geologies, while
targeted enrichment resolves conditional outcome distributions where the
initial realizations indicate mixed behavior.
